\documentclass[10pt]{article}

\usepackage[utf8]{inputenc}
\usepackage[T1]{fontenc}
\usepackage[a4paper,margin=21mm]{geometry}
\usepackage{hyperref}
\usepackage{parskip}
\usepackage{lmodern}
\usepackage{enumitem}

\setlength{\parskip}{4pt}

\hypersetup{
  colorlinks=true,
  urlcolor=black,
  linkcolor=black
}

\begin{document}

\title{List of Changes \\
Manuscript MST-137207 \\
{\normalsize ``Workflow-transfer benchmark of optical surface-roughness measurements against an internal tactile baseline across engineering materials and manufacturing processes''}}

\date{31 July 2026}
\maketitle

\noindent All line numbers refer to the revised \texttt{manuscript.tex}. New or substantially rewritten passages are indicated.

\subsection*{Introduction (Section~1)}

\begin{enumerate}[leftmargin=*, label=\textbf{C\arabic*.}, ref=C\arabic*]

\item \textbf{Profile-vs.-areal rationale (R1\#2).} New paragraph at lines~85--87 explaining why profile parameters were chosen: (i)~production tolerances and interlaboratory comparisons are expressed in profile terms, and (ii)~profile parameters were the only class retained consistently across all four optical families.

\item \textbf{Industrial importance of $R_a$, $R_{sm}$, $R_{z1max}$ (R1\#4).} New paragraph at lines~87--91 describing the distinct engineering functions of the three parameters and linking them to ISO~21920-2/3.

\item \textbf{Structured literature comparison table (R1\#3).} New Table~1 at lines~99--135 with columns for Study, Materials, Parameters, Methods compared, Key finding, and Gap. Includes Marichamy et~al.\ (2025) alongside Garc\'ia et~al., Schmidt et~al., M\'inguez-Mart\'inez et~al., Boedecker et~al., Nouira et~al., Hoefling et~al., Hu et~al., and the present work.

\item \textbf{Research gap with three operational questions (R1\#1).} Substantially expanded paragraph at lines~136--139 articulating the three decisions a practitioner must make when integrating optical roughness measurement into a quality-control workflow, and stating the intended practical output.

\end{enumerate}

\subsection*{Methods (Section~2)}

\begin{enumerate}[leftmargin=*, label=\textbf{C\arabic*.}, ref=C\arabic*]
\setcounter{enumi}{4}

\item \textbf{Data dictionary and independent verification (R1\#15).} New paragraph at lines~207--211 stating that all numerical summaries are generated by documented Python scripts, a data dictionary is provided in the repository README and supplement, and independent verification is possible from the public CSV outputs.

\item \textbf{Band justification with calibration context (R1\#8).} New paragraph at lines~209--213 linking the 10\% threshold to typical calibration-grade agreement and accredited laboratory expanded uncertainty; linking the 30\% threshold to industrial quality-control acceptability; and stating the bands are pragmatic screening tools.

\end{enumerate}

\subsection*{Results (Section~3)}

\begin{enumerate}[leftmargin=*, label=\textbf{C\arabic*.}, ref=C\arabic*]
\setcounter{enumi}{6}

\item \textbf{Best/worst material--process identification (R1\#13, R2\#2).} New paragraph at lines~250--256 identifying best-agreement groups (finish wire-EDM C45, ground C45, honed MO58A, glass-bead-blasted 14301) and worst-agreement groups (rough-face-turned Ti6Al4V, ELLOR graphite, rough-milled 14301). Includes optical constants for Ti6Al4V ($n \approx 3.4$, $k \approx 4.2$) and 14301 ($n \approx 2.5$, $k \approx 3.5$ at 633~nm), with physical-attribution discussion.

\item \textbf{Worst-case case study in main text (R2\#3).} New paragraph at lines~252--260 providing a detailed worst-case illustration using rough-face-turned Ti6Al4V (\texttt{ti6al4v\_tf}), structured identically to the best-case analysis: retrospective lowest-discrepancy, fixed-workflow, and leave-one-surface transfer performance.

\end{enumerate}

\subsection*{Discussion (Section~4)}

\begin{enumerate}[leftmargin=*, label=\textbf{C\arabic*.}, ref=C\arabic*]
\setcounter{enumi}{8}

\item \textbf{$R_{sm}$ three lines of evidence (R1\#12).} Restructured paragraph at lines~348--353 presenting (i)~near-universal 99.9\% discrepancy across all systems/materials/processes, (ii)~substantial reduction under $\times 1000$ rescaling, and (iii)~known sensitivity to ISO~21920-3 profile-element detection rules.

\item \textbf{Unit harmonization recommendation (R1\#9).} Expanded paragraph at lines~350--353 with practical recommendation to verify and document $R_{sm}$ unit convention; where unconfirmed, exclude $R_{sm}$ from optical-substitution claims.

\item \textbf{Unit check for other parameters (R1\#5).} New statement at lines~350--353 confirming that $R_a$, $R_q$, $R_t$, $R_v$, $R_z$, $R_{z1max}$ are reported in micrometres by both tactile and optical instruments, with spot-check verification.

\item \textbf{$R_{sk}$ alternative normalisation (R2\#1).} New paragraph at lines~354--356 proposing absolute difference $\Delta_{sk}=O_{sk}-S_{sk}$ and a rank-correlation-based consistency index as alternatives to percentage normalisation.

\item \textbf{Optical-setting limitations (R1\#6).} New paragraph at lines~356--358 explaining that because objective NA, field of view, lateral sampling, and form-removal settings were not retained uniformly, the benchmark cannot separate principle-level from processing-level discrepancies.

\item \textbf{Optical-repeatability limitations (R1\#7).} New paragraph at lines~358--360 stating that optical repeat measurements were retained less consistently, with a limited sensitivity check indicating optical repeatability SDs are generally smaller than between-workflow discrepancy medians.

\item \textbf{Transfer-penalty patterns (R1\#10).} Expanded paragraph at lines~360--365 detailing which parameters, materials, and processes show better transferability: amplitude parameters transfer better for C45 steel and ground surfaces; $R_a$ and $R_q$ show larger transfer penalties than $R_t$ and $R_z$; cross-material penalties exceed cross-process penalties.

\item \textbf{System-selection guidance (R1\#11).} New paragraph at lines~362--365 discussing the practical selection dilemma when a system performs well for one parameter but poorly for another, with guidance on interferometry-based and confocal configurations.

\item \textbf{ISO~21920-3 tolerance-grade mapping (R1\#14).} New paragraph at lines~366--370 mapping observed discrepancies to ISO~21920-3 tolerance grades TG1--TG6, with caveat that the mapping is illustrative rather than normative.

\item \textbf{Prospective-validation design (R1\#16).} Substantially expanded paragraph at lines~368--377 specifying six requirements: (i)~predefined scope with GPS settings, (ii)~15--20 independent specimens, (iii)~within-instrument repeatability ($n \geq 10$), (iv)~gauge R\&R study, (v)~GUM-compliant uncertainty budget, (vi)~pre-specified acceptance criteria.

\end{enumerate}

\subsection*{Other changes}

\begin{enumerate}[leftmargin=*, label=\textbf{C\arabic*.}, ref=C\arabic*]
\setcounter{enumi}{18}

\item \textbf{Marichamy et~al.\ (2025) reference.} Added to \texttt{bibliography.bib} and cited in the new Table~1 and Introduction text.

\item \textbf{Funding statement.} Added at end of manuscript: Polish Ministry of Science, programme ``Polish Metrology~II'', project no.~PM-II/SP/0090/2024/02.

\item \textbf{Data availability statement.} Added with real GitHub and Zenodo links.

\item \textbf{Conflicts of interest statement.} Added.

\item \textbf{CRediT authorship contribution statement.} Added with contribution percentages and roles for all nine authors.

\end{enumerate}

\end{document}
